% =====================================================================
%  HISTORIAL ACADÉMICO Y PROFESIONAL (HAP) -- Concurso LE-403/729 (UPC)
%  Borrador v3 (05/11/2026), en castellano, elaborado a partir de
%  Short_CV_Carlos_Monteiro-2.pdf.
%  Completar todos los \TODO{} con datos reales y acreditables. Solo se
%  valoran los méritos acreditados a 25/11/2026 (base 3.5).
%  Datos personales: SOLO nombre y apellidos (base 3.5).
%  Antes de presentar: descomentar \finaltrue para ocultar las marcas.
%  Nombre del fichero: Monteiro_Carlos_HAP.pdf
% =====================================================================
\documentclass[11pt]{article}
\usepackage{lector}
%\finaltrue
\setupheader{HAP}

\title{\textbf{Historial Acad\'emico y Profesional}}
\author{\Large \CandName}
\date{}

\begin{document}
\maketitle
\thispagestyle{plain}
\vspace{0.5cm}
\begin{center}\postdata\end{center}

\vspace{0.8cm}
\noindent\textbf{Resumen.}
Soy investigador posdoctoral en la Escuela de Ingenier\'ia de Barcelona Este (EEBE--UPC)
y doctor en Ingenier\'ia Mec\'anica, de Fluidos y Aeron\'autica por la UPC (2026,
Sobresaliente \emph{cum laude}), con la tesis \emph{Pseudo-boiling Wall-Bounded
Turbulence in Supercritical Flows}. Mi investigaci\'on cubre las tres l\'ineas de la plaza:
\begin{itemize}
\item \textbf{Turbulencia multif\'isica en fluidos compresibles y supercr\'iticos:} 6
art\'iculos en revista, 3 de ellos en \emph{Physics of Fluids} (uno seleccionado como
\emph{Editor's Pick} en 2026).
\item \textbf{An\'alisis basado en datos, eventos extremos y causalidad:} estancia en el
CTR Summer Program 2026 de la Universidad de Stanford, sobre inferencia causa--efecto en
eventos extremos de flujo de calor en pared; bases de datos abiertas en
\emph{Scientific Data} y BLASTNet.
\item \textbf{Computaci\'on distribuida y acelerada en GPU:} primer autor del c\'odigo
RHEA-X, residente en GPU y distribuido (\emph{Computer Physics Communications}, 2026), y
participaci\'on en asignaciones competitivas de c\'alculo que suman m\'as de 10 millones de
horas CPU/GPU en MareNostrum~5 y en Canad\'a.
\end{itemize}
En la EEBE imparto pr\'acticas de laboratorio de Mec\'anica de Fluidos desde el curso
2024--25 y, en el curso 2026--27, la teor\'ia completa de Termodin\'amica y Transferencia
de Calor. En primavera de 2027 impartir\'e la teor\'ia de Ingenier\'ia de Fluidos. Aporto
adem\'as experiencia industrial en simulaci\'on t\'ermica (Continental Engineering Services).

\newpage
\tableofcontents
\newpage

% =====================================================================
\section{Formaci\'on y trayectoria}

\subsection{Formaci\'on acad\'emica}
\begin{itemize}
  \item \textbf{2023--2026} -- \textbf{Doctor en Ingenier\'ia Mec\'anica, de Fluidos y
  Aeron\'autica}, Universitat Polit\`ecnica de Catalunya (UPC), Barcelona. Tesis:
  \emph{Pseudo-boiling Wall-Bounded Turbulence in Supercritical Flows}. Director: Dr.\
  Llu\'is Jofre \TODO{confirmar director(es)}. Calificaci\'on: \textbf{Sobresaliente
  \emph{cum laude}}. \certlink{phd}{ver t\'itulo de doctor}
  \item \textbf{2019--2021} -- \textbf{M\'aster en Ingenier\'ia Mec\'anica}, Universidade
  NOVA de Lisboa (Portugal). Trabajo final: \emph{Estudio num\'erico de un sistema de
  reducci\'on de resistencia aerodin\'amica para un veh\'iculo de F\'ormula Student}.
  \certlink{msc}{ver t\'itulo}
  \item \textbf{2019--2020} -- Estancia Erasmus+ en Ingenier\'ia Aeroespacial, Politecnico
  di Milano (Italia). \certlink{erasmus}{ver certificado}
  \item \TODO{T\'itulo de grado/licenciatura: denominaci\'on, universidad, a\~no}
  \certlink{bsc}{ver t\'itulo}
\end{itemize}

\subsection{Trayectoria acad\'emica}
\begin{itemize}
  \item \textbf{2026 -- actualidad}: Investigador posdoctoral en simulaci\'on num\'erica
  (CFD) de flujos turbulentos supercr\'iticos. Departamento de Mec\'anica de Fluidos,
  EEBE--UPC, Barcelona. \TODO{mes de inicio, tipo de contrato y financiaci\'on}
  \certlink{position}{ver certificado}
  \item \textbf{2023 -- 2026}: Doctorando, Departamento de Mec\'anica de Fluidos,
  EEBE--UPC. \TODO{tipo de contrato predoctoral (FPI/FPU/FI-AGAUR/proyecto), si procede}
  \item \textbf{2026}: Investigador visitante, Center for Turbulence Research (CTR)
  Summer Program, Universidad de Stanford (EE.\,UU.). \certlink{ctr}{ver certificado}
\end{itemize}

\subsection{Trayectoria profesional}
\begin{itemize}
  \item \textbf{2021--2022}: Ingeniero de simulaci\'on t\'ermica, Continental Engineering
  Services, Oporto (Portugal). Gesti\'on t\'ermica de componentes el\'ectricos,
  inversores y refrigeraci\'on de motores el\'ectricos. \certlink{continental}{ver certificado}
  \item \textbf{2020--2021}: Ingeniero de aerodin\'amica y CFD, equipo de F\'ormula Student
  FST Lisboa (Portugal). Aerodin\'amica externa y sistemas de reducci\'on de resistencia.
  \item \textbf{2019}: Ingeniero de dise\~no CAD 3D en pr\'acticas, Cubotonic Lda.,
  Malveira (Portugal).
\end{itemize}

\subsection{Acreditaciones}
\begin{itemize}
  \item \TODO{Informe favorable de AQU (Lector) o ANECA (Ayudante Doctor), si se dispone
  (m\'erito preferente, +10 puntos). En caso contrario, eliminar este apartado.}
\end{itemize}

% =====================================================================
\section{Actividad docente}
\NOTE{Criterio 1 (hasta 10 puntos). Adjuntar el certificado oficial de docencia de la UPC
con las horas por asignatura y cuatrimestre. Las horas del curso 2026--27 deben estar
certificadas por el departamento.}

\subsection{Docencia en titulaciones oficiales (EEBE--UPC)}
\begin{center}\small
\begin{tabularx}{\textwidth}{@{}p{3.6cm}p{3.0cm}Xp{1.3cm}@{}}
\toprule
\textbf{Asignatura} & \textbf{Actividad} & \textbf{Curso / cuatrimestre} & \textbf{Horas} \\
\midrule
Termodin\'amica y Transferencia de Calor & Teor\'ia (asignatura completa) & 2026--27 (en curso) \TODO{cuatrimestre, grupo(s), grado(s)} & \TODO{} \\
Mec\'anica de Fluidos & Pr\'acticas de laboratorio & 2026--27 (en curso) & \TODO{} \\
Fundamentos de Ingenier\'ia T\'ermica \TODO{comprobar denominaci\'on oficial} & Pr\'acticas de laboratorio & 2026--27 (en curso) & \TODO{} \\
Ingenier\'ia de Fluidos & Teor\'ia (asignada) & 2026--27, cuatrimestre de primavera & \TODO{} \\
Mec\'anica de Fluidos & Pr\'acticas de laboratorio & 2025--26, Q1 y Q2 & 20 + 20 \\
Mec\'anica de Fluidos & Pr\'acticas de laboratorio & 2024--25, Q2 & 20 \\
\midrule
\multicolumn{3}{@{}r}{\textbf{Total impartido a 25/11/2026 (sin contar la docencia asignada)}} & \TODO{} \\
\bottomrule
\end{tabularx}
\end{center}
\certlink{teaching}{ver certificado de docencia} \quad
\certlink{assignment}{ver encargo docente 2026--27}

\subsection{Direcci\'on de trabajos acad\'emicos}
\begin{itemize}
  \item \textbf{2023--2025}: codirecci\'on de Trabajos de Fin de Grado (TFG) en la UPC.
  \TODO{para cada uno: estudiante, t\'itulo, grado, a\~no, calificaci\'on}
  \certlink{tfg}{ver certificado}
\end{itemize}

\subsection{Material docente}
\begin{itemize}
  \item \TODO{material elaborado para la teor\'ia de TTC (transparencias, colecciones de
  problemas, ex\'amenes), guiones de pr\'acticas de MF, cuestionarios en Atenea, cuadernos de Python}
\end{itemize}

\subsection{Indicadores de calidad docente}
\begin{itemize}
  \item \TODO{resultados de las encuestas al estudiantado (laboratorios de MF 2024--25 y
  2025--26), si se dispone de ellos} \certlink{surveys}{ver encuestas}
\end{itemize}

% =====================================================================
\section{Formaci\'on continua}
\subsection{Formaci\'on docente}
\NOTE{Criterio 3 (hasta 5 puntos). Solo cuentan los cursos finalizados y certificados a 25/11/2026.
El certificado del ICE exige una valoraci\'on positiva del trabajo propuesto: entregarlo y
pedir el certificado en cuanto acabe el curso. Un solo curso de 3 h punt\'ua poco; conviene
a\~nadir otros cursos del ICE que terminen antes del 25/11.}
\begin{itemize}
  \item 2026: \emph{M2.4 -- Atenea: curso de nivel medio} (1.\textsuperscript{a} edici\'on).
  Instituto de Ciencias de la Educaci\'on (ICE), UPC. 3 h, en l\'inea (9--11/11/2026).
  Uso avanzado del campus virtual Atenea (Moodle): foros, consultas, control de asistencia,
  lecciones, grupos y agrupamientos, actividades condicionales y excepciones de tiempo, aplicado
  a la organizaci\'on de las asignaturas de la EEBE. \certlink{ice1}{ver certificado}
\end{itemize}

\subsection{Formaci\'on cient\'ifica y t\'ecnica}
\begin{itemize}
  \item \textbf{2024}: Compute Ontario Summer School (Digital Research Alliance of
  Canada): computaci\'on avanzada para la investigaci\'on, HPC y software de investigaci\'on.
  \item \textbf{\TODO{a\~no}}: CFDParSchool, Gran Sasso Science Institute, L'Aquila (Italia):
  paralelizaci\'on de la din\'amica de fluidos num\'erica y HPC.
  \item \textbf{\TODO{a\~no}}: \emph{Hands-on Machine Learning for Fluid Dynamics}, von
  Karman Institute for Fluid Dynamics (en l\'inea).
  \item \textbf{2022}: \emph{Machine Learning for Mechanical Engineers} (24 h)
  \TODO{comprobar a\~no} y \emph{Python for Engineers} (20 h), Flowthermolab.
  \item \textbf{2022}: Ansys Turbulence Modelling Series: \emph{Beyond Basics,
  Scale-Resolving Simulations} (12 h).
  \item \textbf{2019--2021}: Ansys FEA Workshop (8 h, ESSS); FSG Academy -- Aerodin\'amica
  (Formula Student Germany); conferencia de aerodin\'amica de F1 (Red Bull, Politecnico di
  Milano); aerodin\'amica de veh\'iculos y ensayos en t\'unel de viento (Politecnico di Milano).
\end{itemize}

% =====================================================================
\section{Actividad investigadora}
\NOTE{Criterio 2 (hasta 10 puntos).}

\subsection{Identificadores e indicadores}
ORCID: \TODO{} \quad Scopus Author ID: \TODO{} \quad Google Scholar: \TODO{}\\
7 art\'iculos en revistas indexadas (2024--2026; 4 como primer autor y los 7 como primer o
segundo autor), 2 comunicaciones en congresos, 2 bases de datos abiertas y 1
\emph{Editor's Pick} (\emph{Physics of Fluids}, 2026).
\TODO{citas, \'indice h, cuartiles (JCR/SJR) por revista}

\subsection{L\'ineas de investigaci\'on}
Mi investigaci\'on abarca las tres l\'ineas de investigaci\'on de la plaza:
\begin{enumerate}
  \item \textbf{Flujos turbulentos multif\'isicos; fluidos compresibles y supercr\'iticos.}
  Simulaci\'on num\'erica directa (DNS) de turbulencia de pared en fluidos transcr\'iticos y
  supercr\'iticos a alta presi\'on: pseudoebullici\'on (\emph{pseudo-boiling}),
  termodin\'amica de gas real, flujos secundarios inducidos por flotabilidad en conductos,
  flujos microconfinados y su transici\'on a la turbulencia, unidades m\'inimas de flujo y
  criterios de resoluci\'on para DNS [J2, J3, J4, J5, J6, J7, C1; tesis doctoral].
  \item \textbf{Predicci\'on basada en datos de flujos de ingenier\'ia; eventos extremos;
  causalidad.} Inferencia causa--efecto en eventos extremos de flujo de calor en pared en
  flujos supercr\'iticos (CTR Summer Program 2026, Stanford). Intermitencia lagrangiana y
  eventos termofluidodin\'amicos extremos (asignaci\'on RES IM-2026-1-0034).
  Caracterizaci\'on de escalas espacio-temporales en experimentos mediante t\'ecnicas basadas
  en datos [J7]. Bases de datos abiertas para aprendizaje autom\'atico [J3, D1, D2] y
  reconstrucci\'on de flujos a partir de sensores dispersos (RES IM-2025-1/2-0017).
  \item \textbf{Modelizaci\'on y computaci\'on; programaci\'on distribuida (CPU) y acelerada
  (GPU).} Desarrollo del c\'odigo abierto RHEA y de su versi\'on distribuida y residente en
  GPU, RHEA-X (MPI + OpenACC/CUDA) [J1]. Portabilidad a GPU en los MareNostrum Hackathons
  del BSC (2024, 2026) y en el Open Accelerated Computing Summit (2023). Simulaciones a gran
  escala en MareNostrum~5, Finisterrae~III y los sistemas nacionales de Canad\'a.
\end{enumerate}

\subsection{Grupos de investigaci\'on}
\begin{itemize}
  \item \TODO{grupo de investigaci\'on del Departamento de Mec\'anica de Fluidos (nombre,
  c\'odigo SGR si est\'a reconocido), fechas de pertenencia}
\end{itemize}

\subsection{Publicaciones}
\subsubsection*{Art\'iculos en revistas indexadas}
\begin{enumerate}[label={[J\arabic*]}]
  \item \textbf{C.\ Monteiro}, X.\ \'Alvarez-Farr\'e, D.\ Mart\'in, F.\ Pizzi, L.\ Jofre.
  RHEA-X: Distributed GPU-resident high-fidelity flow solver for multiphysics turbulence.
  \emph{Computer Physics Communications}, 110233, 2026.
  \href{https://doi.org/10.1016/j.cpc.2026.110233}{doi:10.1016/j.cpc.2026.110233}. \TODO{cuartil} \emph{(L\'inea 3)}
  \item \textbf{C.\ Monteiro}, G.\ Barea, L.\ Jofre. Buoyancy-driven secondary flows and
  momentum redistribution in high-pressure transcritical duct turbulence.
  \emph{Physics of Fluids}, 065127, 2026.
  \href{https://doi.org/10.1063/5.0337547}{doi:10.1063/5.0337547}. \textbf{Editor's Pick.} \emph{(L\'inea 1)}
  \item A.\ Abdellatif, \textbf{C.\ Monteiro}, M.\ Bernades, L.\ Jofre. Microconfined
  high-pressure transcritical channel flow database: laminar, transitional \& turbulent
  regimes. \emph{Scientific Data}, 11(1), 2024.
  \href{https://doi.org/10.1038/s41597-024-03757-4}{doi:10.1038/s41597-024-03757-4}. \emph{(L\'ineas 1 y 2)}
  \item \textbf{C.\ Monteiro}, L.\ Jofre. Flow regime analysis of high-pressure
  transcritical fluids in microducts. \emph{International Journal of Heat and Mass
  Transfer}, 224, 125295, 2024.
  \href{https://doi.org/10.1016/j.ijheatmasstransfer.2024.125295}{doi:10.1016/j.ijheatmasstransfer.2024.125295}. \emph{(L\'inea 1)}
  \item \textbf{C.\ Monteiro}, L.\ Jofre. Resolution standards for direct numerical
  simulation of wall turbulence in high-pressure transcritical fluids. \emph{Physics of
  Fluids}, 36(12), 2024.
  \href{https://doi.org/10.1063/5.0244472}{doi:10.1063/5.0244472}. \emph{(L\'ineas 1 y 3)}
  \item R.\ El Mansy, \textbf{C.\ Monteiro}, F.\ Mellibovsky, L.\ Jofre. Minimal flow unit
  of wall-bounded high-pressure transcritical turbulence. \emph{Physics of Fluids},
  36(12), 2024.
  \href{https://doi.org/10.1063/5.0243832}{doi:10.1063/5.0243832}. \emph{(L\'inea 1)}
  \item E.\ Hurt\'an, \textbf{C.\ Monteiro}, M.\ Jofre, J.\ Casals-Terr\'e, L.\ Jofre.
  Data-informed characterization of spatio-temporal scales in experiments of microconfined
  high-pressure transcritical turbulence. \emph{Experimental Thermal and Fluid Science},
  159, 111282, 2024.
  \href{https://doi.org/10.1016/j.expthermflusci.2024.111282}{doi:10.1016/j.expthermflusci.2024.111282}. \emph{(L\'ineas 1 y 2)}
\end{enumerate}

\subsubsection*{Art\'iculos en preparaci\'on}
\begin{enumerate}[label={[R\arabic*]}]
  \item \TODO{trabajo del CTR 2026: inferencia causa--efecto en eventos extremos de flujo de
  calor en pared en flujos supercr\'iticos (CTR Proceedings of the Summer Program 2026 /
  art\'iculo en preparaci\'on)} \emph{(L\'inea 2)}
\end{enumerate}

\subsubsection*{Comunicaciones en congresos}
\begin{enumerate}[label={[C\arabic*]}]
  \item \textbf{C.\ Monteiro}, L.\ Jofre. Buoyant high-pressure transcritical fluid duct
  flows. \emph{9th European Congress on Computational Methods in Applied Sciences and
  Engineering (ECCOMAS 2024)}, Lisboa, 2024.
  \href{https://doi.org/10.23967/eccomas.2024.054}{doi:10.23967/eccomas.2024.054}. Comunicaci\'on oral.
  \item \textbf{C.\ Monteiro}, M.\ Brito, D.\ Vieira. Numerical iterative analysis of drag
  reduction system for a racing car rear wing. \emph{SAE Technical Paper Series}, 2023.
  \href{https://doi.org/10.4271/2023-01-5079}{doi:10.4271/2023-01-5079}.
\end{enumerate}

\subsubsection*{Bases de datos}
\begin{enumerate}[label={[D\arabic*]}]
  \item W.\ T.\ Chung, M.\ Ihme, K.\ S.\ Jung, \emph{et al.} (incl.\ \textbf{C.\ Monteiro}).
  BLASTNet multiphysics simulation dataset, v3.0. Zenodo, 2026.
  \href{https://doi.org/10.5281/zenodo.21348075}{doi:10.5281/zenodo.21348075}. \emph{(L\'ineas 1 y 2)}
  \item A.\ Abdellatif, \textbf{C.\ Monteiro}, L.\ Jofre. Microconfined high-pressure
  transcritical channel flow dataset, v3. Figshare+, 2023.
  \href{https://doi.org/10.25452/figshare.plus.24486904.v3}{doi:10.25452/figshare.plus.24486904.v3}.
\end{enumerate}

\subsection{Conferencias invitadas y presentaciones}
\begin{itemize}
  \item \textbf{2026}: Seminario invitado, Seminars@DEM, Instituto Superior T\'ecnico,
  Lisboa: \emph{Pseudo-boiling wall-bounded turbulence in supercritical flows}.
  \item \textbf{2026}: CTR Summer Program, Universidad de Stanford: \emph{Cause--effect
  inference in extreme wall heat-flux events in supercritical flows}. \TODO{¿presentaci\'on final?}
  \item \textbf{2024}: ECCOMAS 2024, Lisboa: \emph{Buoyant high-pressure transcritical fluid duct flows}.
  \item \textbf{2024}: Seminarios de Mec\'anica Computacional, EEBE--UPC: \emph{Turbulent
  square-duct flows at high-pressure transcritical conditions}.
  \item \textbf{2023}: Open Accelerated Computing Summit (en l\'inea): \emph{GPU acceleration
  of the massively-parallel flow solver RHEA with MPI+OpenACC}.
\end{itemize}

\subsection{Asignaciones competitivas de recursos de supercomputaci\'on}
\begin{center}\small
\begin{tabularx}{\textwidth}{@{}p{1.4cm}p{3.3cm}Xp{2.4cm}@{}}
\toprule
\textbf{A\~no} & \textbf{Convocatoria / sistema} & \textbf{Proyecto} & \textbf{Recursos} \\
\midrule
2026--27 & RES IM-2026-1-0034, MareNostrum 5 (BSC) & Extreme thermofluid events in transcritical flows: Lagrangian intermittency in supercritical fluid turbulence (IP: L.\ Jofre) & 4 M h-n\'ucleo \\
2026 & Digital Research Alliance of Canada, RAC-2026 kfw-162-ab & GPU-accelerated simulations of microplastics transport in ocean turbulence (IP: M.\ Rahmani) & 1340 a\~nos-n\'ucleo \\
2025--26 & RES IM-2025-2-0017, MareNostrum 5 (BSC) & Dataset generation for turbulent flow reconstruction in external aerodynamics using sparse sensor data (IP: J.\ Calafell) & 1,7 M h-GPU \\
2025 & RES IM-2025-1-0017, MareNostrum 5 (BSC) & Mismo proyecto (IP: J.\ Calafell) & 4 M h-GPU \\
2023--24 & RES IM-2023-2-0005, CTE-AMD (BSC) & Direct numerical simulations of microconfined high-pressure transcritical fluid turbulence (IP: L.\ Jofre) & 490.000 h-n\'ucleo \\
2023 & RES IM-2023-1-0005, Finisterrae III-GPU (CESGA) & Mismo proyecto (IP: L.\ Jofre) & 500.000 h-n\'ucleo \\
\bottomrule
\end{tabularx}
\end{center}
\TODO{indicar el papel en cada una (miembro del equipo / cosolicitante) y adjuntar las resoluciones}

\subsection{Software de investigaci\'on}
\begin{itemize}
  \item \textbf{RHEA / RHEA-X}: c\'odigo abierto, MPI + OpenACC/CUDA, residente en GPU, para
  la simulaci\'on de alta fidelidad de turbulencia multif\'isica compresible y supercr\'itica
  [J1]. Papel: desarrollador principal de RHEA-X \TODO{confirmar papel; a\~nadir URL del repositorio}.
  \item Portabilidad a GPU y optimizaci\'on a nivel de \emph{kernel} de un modelo sustituto
  termodin\'amico (8.\textsuperscript{o} MareNostrum Hackathon, BSC, 2026). Aceleraci\'on de
  RHEA con OpenACC (6.\textsuperscript{o} MareNostrum Hackathon, BSC, 2024).
\end{itemize}

\subsection{Proyectos de investigaci\'on}
\begin{itemize}
  \item \TODO{proyectos nacionales o europeos en los que participa como miembro del equipo
  (PID, ERC, etc.): t\'itulo, referencia, IP, periodo}
\end{itemize}

\subsection{Premios y reconocimientos}
\begin{itemize}
  \item \textbf{2026}: \emph{Editor's Pick} de \emph{Physics of Fluids} para [J2].
  \item Tesis doctoral calificada con Sobresaliente \emph{cum laude} (2026).
\end{itemize}

% =====================================================================
\section{Internacionalizaci\'on}
\NOTE{Criterio 4 (hasta 5 puntos).}
\subsection{Estancias de investigaci\'on internacionales}
\begin{itemize}
  \item \textbf{2026 (1 mes)} -- \textbf{Investigador visitante, Center for Turbulence
  Research Summer Program, Universidad de Stanford (EE.\,UU.).} Proyecto: \emph{Cause--effect
  inference in extreme wall heat-flux events in supercritical flows}, sobre fluidos
  supercr\'iticos, eventos extremos y causalidad. \TODO{investigadores anfitriones/colaboradores}
  \certlink{ctr}{ver certificado}
  \item \textbf{Junio--agosto de 2027 (3 meses, financiada y ya organizada)} --
  \textbf{University of British Columbia (Canad\'a).} Turbulencia multif\'isica y causalidad.
  Anfitri\'on: \TODO{p.\,ej.\ Prof.\ Mona Rahmani -- confirmar}. Financiaci\'on:
  \TODO{programa}. \certlink{ubc}{ver carta de invitaci\'on y de financiaci\'on}
\end{itemize}

\subsection{Formaci\'on y movilidad internacional}
\begin{itemize}
  \item 2019--2020: estancia Erasmus+ en el Politecnico di Milano (Italia).
  \item M\'aster en la Universidade NOVA de Lisboa (Portugal) y doctorado en la UPC.
  \item 2024: Compute Ontario Summer School (Canad\'a). CFDParSchool en el GSSI (L'Aquila,
  Italia). Curso de aprendizaje autom\'atico del von Karman Institute (B\'elgica).
\end{itemize}

\subsection{Colaboraciones internacionales}
\begin{itemize}
  \item Universidad de Stanford (CTR; base de datos BLASTNet con el grupo de M.\ Ihme) [D1].
  \item University of British Columbia / Digital Research Alliance of Canada (asignaci\'on RAC-2026).
  \item Instituto Superior T\'ecnico, Lisboa (seminario invitado, 2026).
  \item Previstas para 2027--2029: estancias en grupos europeos, p.\,ej.\ el del Prof.\ Mario
  Di Renzo (Universit\`a del Salento, Italia) \TODO{confirmar}; v\'ease el Plan de Trabajo.
\end{itemize}

% =====================================================================
\section{Otros m\'eritos}
\NOTE{Criterio 5 (hasta 5 puntos).}
\subsection{Idiomas}
\begin{itemize}
  \item \textbf{Portugu\'es}: lengua materna.
  \item \textbf{Ingl\'es} y \textbf{castellano}: nivel fluido. \TODO{certificados, si se dispone}
  \item \textbf{Catal\'an}: nivel B1 \TODO{confirmar nivel, organismo emisor y fecha}.
  Cursos: \TODO{curso, nivel, horas, a\~no}. Formaci\'on en curso para alcanzar el nivel C1.
  \certlink{catalan}{ver certificados}
\end{itemize}
\subsection{Experiencia industrial y transferencia de conocimiento}
\begin{itemize}
  \item Simulaci\'on t\'ermica de componentes el\'ectricos, inversores y refrigeraci\'on de
  motores el\'ectricos en Continental Engineering Services (2021--2022), experiencia
  directamente aplicable a Ingenier\'ia T\'ermica y a TTC.
  \item Aerodin\'amica de autom\'ocion y sistemas de reducci\'on de resistencia (FST Lisboa,
  publicaci\'on SAE [C2]).
\end{itemize}
\subsection{Competencias t\'ecnicas}
C++, Python, MATLAB, CUDA, OpenACC, MPI; entornos HPC Linux/Slurm; STAR-CCM+, Ansys;
SolidWorks, Siemens NX, Creo.

% =====================================================================
\section{Descripci\'on de la trayectoria acad\'emica y profesional}
Mi trayectoria ha evolucionado de la ingenier\'ia aplicada a la investigaci\'on fundamental
en termofluidos, siempre con la simulaci\'on num\'erica como eje. Tras cursar el M\'aster en
Ingenier\'ia Mec\'anica en Lisboa, con un a\~no Erasmus+ de Ingenier\'ia Aeroespacial en
Mil\'an, trabaj\'e en aerodin\'amica en F\'ormula Student. Despu\'es me incorpor\'e a
Continental Engineering Services como ingeniero de simulaci\'on t\'ermica, donde model\'e la
refrigeraci\'on de electr\'onica de potencia y de motores el\'ectricos. Esa experiencia
industrial me mostr\'o hasta qu\'e punto la transferencia de calor y la mec\'anica de fluidos
condicionan el dise\~no real, y me llev\'o a querer comprender estos fen\'omenos en
profundidad.

En 2023 inici\'e el doctorado en el Departamento de Mec\'anica de Fluidos de la UPC, sobre
la turbulencia de pared en fluidos supercr\'iticos. Cerca del punto pseudocr\'itico, estos
flujos combinan turbulencia, fuertes efectos termodin\'amicos de gas real y transferencia
de calor. Durante el doctorado publiqu\'e seis art\'iculos en revista sobre este tema: reg\'imenes
de flujo en microconductos, criterios de resoluci\'on para DNS, unidades m\'inimas de flujo y
flujos secundarios inducidos por flotabilidad en conductos; este \'ultimo fue seleccionado
como \emph{Editor's Pick} en \emph{Physics of Fluids}. Para calcular estos flujos, lider\'e el
desarrollo de RHEA-X, un c\'odigo distribuido y residente en GPU publicado en
\emph{Computer Physics Communications}, y particip\'e en varias asignaciones competitivas de
supercomputaci\'on. Defend\'i la tesis en 2026 con la calificaci\'on de Sobresaliente
\emph{cum laude}.

En 2026 fui seleccionado para el CTR Summer Program de la Universidad de Stanford, donde
trabaj\'e en la inferencia causa--efecto de eventos extremos de flujo de calor en pared en
flujos supercr\'iticos. Este trabajo conecta mi experiencia con el an\'alisis basado en
datos, los eventos extremos y la causalidad, y lo continuar\'e en la University of British
Columbia en el verano de 2027.

En paralelo, imparto docencia en la EEBE desde 2024: primero pr\'acticas de laboratorio de
Mec\'anica de Fluidos y, en el curso 2026--27, la teor\'ia completa de Termodin\'amica y
Transferencia de Calor. En primavera de 2027 impartir\'e la teor\'ia de Ingenier\'ia de
Fluidos. La plaza de Profesor Lector me permitir\'ia consolidar esta trayectoria docente e
investigadora en el departamento donde comenz\'o.
\TODO{revisar y adaptar a su propia voz}

% =====================================================================
\section{Adecuaci\'on al perfil de la plaza}
\begin{center}\small
\begin{tabularx}{\textwidth}{@{}p{5.0cm}X@{}}
\toprule
\textbf{Perfil de la convocatoria} & \textbf{M\'eritos que lo acreditan} \\
\midrule
\multicolumn{2}{@{}l}{\emph{Docencia}}\\
Termodin\'amica y Transferencia de Calor; Termodin\'amica & Teor\'ia completa de TTC en la EEBE (2026--27) \\
Mec\'anica de Fluidos & Pr\'acticas de laboratorio en la EEBE cada curso desde 2024--25 \\
Ingenier\'ia de Fluidos & Teor\'ia, asignada para la primavera de 2027 \\
Ingenier\'ia T\'ermica; Centrales El\'ectricas y Energ\'ias Renovables; Generaci\'on Termofluidodin\'amica & Pr\'acticas de Fundamentos de Ingenier\'ia T\'ermica (2026--27); simulaci\'on t\'ermica industrial en Continental; investigaci\'on en fluidos supercr\'iticos (ciclos de sCO$_2$) \\
Mec\'anica de Fluidos Computacional (M\'aster Universitario en Tecnolog\'ias Mec\'anicas) & Experiencia en CFD y HPC: RHEA-X [J1], hackathons de GPU, escuelas de HPC \\
\midrule
\multicolumn{2}{@{}l}{\emph{Investigaci\'on}}\\
Flujos turbulentos multif\'isicos; fluidos compresibles y supercr\'iticos & Tesis doctoral; [J2--J7], [C1], [D1, D2]; \emph{Editor's Pick} \\
Predicci\'on basada en datos; eventos extremos; causalidad & CTR-Stanford 2026; RES IM-2026-1-0034; [J3], [J7], [D1]; UBC 2027 \\
Modelizaci\'on y computaci\'on; programaci\'on en CPU y GPU & RHEA-X [J1], [J5]; 6 asignaciones de HPC ($>$10 M h CPU/GPU); OpenACC/CUDA/MPI \\
Direcci\'on de tesis; direcci\'on de grupos y proyectos & Codirecci\'on de TFG (2023--2025); tutorizaci\'on de estudiantes del grupo \TODO{} \\
\bottomrule
\end{tabularx}
\end{center}

\noindent En resumen:
\begin{itemize}
\item Mi investigaci\'on cubre \emph{las tres} l\'ineas de investigaci\'on de la plaza, con
art\'iculos como primer autor en cada una, y contin\'ua a trav\'es del CTR-Stanford (2026) y
la UBC (2027).
\item Ya imparto asignaturas troncales del perfil en la EEBE (teor\'ia de TTC, laboratorios
de MF e Ingenier\'ia de Fluidos a partir de la primavera de 2027) y conozco a su
estudiantado y sus laboratorios.
\item Mi experiencia industrial en simulaci\'on t\'ermica aporta una visi\'on aplicada a las
asignaturas de ingenier\'ia t\'ermica.
\item Mi experiencia en HPC y GPU me permite contribuir a la asignatura de M\'aster de
Mec\'anica de Fluidos Computacional y sostener una l\'inea de investigaci\'on competitiva a
nivel internacional.
\end{itemize}

% =====================================================================
\newpage
\section{Documentaci\'on justificativa}
Los documentos justificativos se adjuntan a continuaci\'on, en el orden en que se citan.

\certpdf{phd}{PhD_title.pdf}
\certpdf{msc}{MSc_title.pdf}
\certpdf{erasmus}{Erasmus_certificate.pdf}
\certpdf{bsc}{BSc_title.pdf}
\certpdf{position}{Postdoc_contract_certificate.pdf}
\certpdf{continental}{Continental_certificate.pdf}
\certpdf{teaching}{UPC_teaching_certificate.pdf}
\certpdf{assignment}{Teaching_assignment_2026-27.pdf}
\certpdf{tfg}{TFG_cosupervision.pdf}
\certpdf{surveys}{Student_surveys.pdf}
\certpdf{ice1}{ICE_M2-4_Atenea_2026.pdf}
\certpdf{ctr}{CTR_Summer_Program_2026_certificate.pdf}
\certpdf{ubc}{UBC_2027_invitation_and_funding.pdf}
\certpdf{catalan}{Catalan_certificates.pdf}

\end{document}
